\section*{Capítulo \ref{ch:b3:dna-repair} --- Estabilidade do genoma: dano, reparo e recombinação do DNA}

\begin{solution}{exo:b3:dna-repair:1}
Uracila: espontânea (\omterm{def:b3:dna-repair:lesions}{desaminação} de C), \omterm{def:b3:dna-repair:excision}{reparo por excisão de bases}
pela uracila-DNA-glicosilase. Dímero de timina: ambiental (UV), \omterm{def:b3:dna-repair:excision}{reparo
por excisão de nucleotídeos}. \omterm{def:b3:dna-repair:lesions}{8-oxoguanina}: espontânea (oxidação),
\omterm{def:b3:dna-repair:excision}{reparo por excisão de bases} pela \omterm{def:b3:dna-repair:excision}{glicosilase} da 8-oxoG. Mau pareamento
G--T atrás da forquilha: erro de replicação, \omterm{prop:b3:dna-repair:fidelity}{reparo de mau pareamento}.
\omterm{def:b3:dna-repair:lesions}{Quebra de fita dupla} em G1: ambiental (radiação) ou espontânea; \omterm{def:b3:dna-repair:dsb}{junção
de extremidades não homólogas}, já que não há cromátide irmã
disponível.
\end{solution}

\begin{solution}{exo:b3:dna-repair:2}
A excisão de nucleotídeos reconhece uma propriedade física comum a
todas as lesões volumosas --- a distorção da hélice, ou uma
RNA-polimerase parada --- e remove um remendo que contém o que quer
que a tenha causado. As lesões pequenas (uracila, 8-oxoG, bases
alquiladas) quase não distorcem a hélice, e por isso precisam ser
reconhecidas por sua química, uma \omterm{def:b3:dna-repair:excision}{glicosilase} por tipo de base errada.
\end{solution}

\begin{solution}{exo:b3:dna-repair:3}
Um mau pareamento é feito de duas bases normais, e só a que está na
fita nova está errada; o sistema precisa saber qual fita é a nova.
\emph{E.~coli} marca a fita parental com grupos metila nos GATC, que a
fita nova não tem por alguns minutos; a MutH corta a fita não metilada.
Num mutante \emph{dam} nenhuma das fitas é metilada: a MutH corta
qualquer uma, o reparo remove a base certa tantas vezes quanto a
errada --- fixando metade dos erros como mutações (uma mutadora) --- e
cortes nas duas fitas em sítios GATC próximos produzem \omterm{def:b3:dna-repair:lesions}{quebras de fita
dupla}.
\end{solution}

\begin{solution}{exo:b3:dna-repair:4}
Cerca de $35\times 2 = 70$ \omterm{def:b3:dna-repair:lesions}{quebras de fita dupla}. Após uma hora a
maioria já foi unida: com meia-vida perto de \qty{30}{min}, cerca de
$70/4 \approx 18$ focos. Após seis horas quase todas foram reparadas
($70/2^{12} < 1$), restando apenas algumas quebras complexas e
persistentes.
\end{solution}

\begin{solution}{exo:b3:dna-repair:5}
(a) $10^{-5}\times 10^{-2}\times 10^{-3} = 10^{-10}$; $6.4\times
10^{9}\times 10^{-10} = 0.64$ mutações por divisão. (b) Sem \omterm{prop:b3:dna-repair:fidelity}{reparo de
mau pareamento}, $10^{-7}$; $640$ por divisão. (c) Sem revisão,
$10^{-8}$; $64$ por divisão. A mutante de \omterm{prop:b3:dna-repair:fidelity}{reparo de mau pareamento} é a
mais perigosa, mil vezes acima do normal.
\end{solution}

\begin{solution}{exo:b3:dna-repair:6}
$k = \ln 2/\qty{0.25}{h} = \qty{2.8}{h^{-1}}$; $D = 10^{4}/24 =
\qty{420}{h^{-1}}$; $L^{*} = 420/2.8 \approx 150$ sítios abásicos
presentes a cada momento. Com o reparo dez vezes mais lento:
$L^{*} \approx \num{1500}$. A forquilha passa uma vez por cada loco e
encontra as lesões presentes ali naquele momento; na média do genoma
ela encontra cerca de $L^{*}$ delas --- uns $150$ normalmente,
$\num{1500}$ com o fármaco.
\end{solution}

\begin{solution}{exo:b3:dna-repair:7}
$k_{\text{NHEJ}} = \ln 2/0.5 = \qty{1.39}{h^{-1}}$, $k_{\text{HR}} =
\ln 2/4 = \qty{0.17}{h^{-1}}$; fração da HR $0.17/1.56 = 0.11$;
meia-vida $\ln 2/1.56 = \qty{27}{min}$. Uma forquilha colapsada dá uma
quebra de \emph{uma só extremidade}: não há segunda extremidade para a
NHEJ unir, de modo que só a recombinação com a irmã pode restaurar a
forquilha --- e ela é, de todo modo, a única via acurada.
\end{solution}

\begin{solution}{exo:b3:dna-repair:8}
(a) Lynch: microssatélites instáveis (comprimentos diferentes entre as
células do tumor), cânceres colorretal, de endométrio, gástrico e de
ovário, resposta normal à luz do sol. (b) XP-A: microssatélites
estáveis, queimadura extrema com exposição mínima, cânceres de pele e
de olho na infância e, neste grupo, degeneração neurológica
progressiva. (c) XP-V: microssatélites estáveis, reparo por excisão
normal, de modo que a queimadura é mais branda, mas os dímeros que
escapam da excisão são contornados por polimerases propensas a erro
--- cânceres de pele, de início mais tardio que no grupo~A.
\end{solution}

\begin{solution}{exo:b3:dna-repair:9}
Células hepáticas: a Cre excisou o éxon~3 dos dois alelos (como um
círculo que se perde), de modo que o gene está inativo em cada
hepatócito desde o nascimento --- um nocaute específico do fígado.
Neurônios: sem Cre, os alelos floxeados estão intactos e funcionais.
Orientações opostas: a Cre inverte o éxon e o desinverte
indefinidamente, de modo que a qualquer momento cerca de metade das
células carrega o éxon invertido e não funcional --- um mosaico, e não
uma deleção limpa.
\end{solution}

\begin{solution}{exo:b3:dna-repair:10}
À medida que a LexA é clivada, sua concentração cai gradualmente. Os
operadores que ligam a LexA fracamente ficam vagos já com uma queda
pequena --- \emph{uvrA}, \emph{uvrB}, \emph{recA} ---, e por isso o
reparo acurado começa primeiro; os operadores de \emph{sulA} e das
polimerases propensas a erro ligam a LexA com força e só ficam vagos
quando ela quase desapareceu, isto é, quando o dano persistiu por
muito tempo. Uma bactéria sem as polimerases propensas a erro morreria
em forquilhas bloqueadas por lesões que não conseguiria contornar e,
sem mutagênese induzida, evoluiria resistência ao antibiótico bem mais
devagar --- uma desvantagem para a bactéria, uma vantagem para o
paciente; inibidores da \omterm{prop:b3:dna-repair:sos}{resposta SOS} são estudados por essa razão.
\end{solution}

\begin{solution}{exo:b3:dna-repair:11}
Cadeia: PARP inibida $\to$ as quebras de fita simples persistem $\to$
uma forquilha de replicação converte cada uma numa \omterm{def:b3:dna-repair:lesions}{quebra de fita
dupla} de uma só extremidade $\to$ o reparo exige \omterm{def:b3:dna-repair:dsb}{recombinação homóloga}
$\to$ a célula tumoral, sem \omterm{def:b3:dna-repair:dsb}{BRCA}, não consegue $\to$ união errada,
perda de cromossomos, morte. As células normais, com um alelo
funcional, recombinam e sobrevivem. Resistência: (1) uma segunda
mutação em \emph{\omterm{def:b3:dna-repair:dsb}{BRCA}} que restaura a fase de leitura e uma proteína
funcional --- detecte sequenciando o tumor ou o DNA tumoral circulante
no sangue, e pelo retorno dos focos de \omterm{def:b3:dna-repair:dsb}{Rad51} após a irradiação; (2) a
perda das proteínas que bloqueiam a ressecção das extremidades
(53BP1--shieldina), que restaura em parte a recombinação sem a BRCA1
--- detecte pela ausência delas na coloração e pelos focos de \omterm{def:b3:dna-repair:dsb}{Rad51}
restaurados; também mutações de PARP1 que impeçam seu aprisionamento
no DNA, ou bombas de efluxo de fármacos, detectadas por sequenciamento
e por expressão, respectivamente.
\end{solution}

\begin{solution}{exo:b3:dna-repair:12}
Dois dúplexes homólogos, um carregando $A$ e o outro $a$, trocam fitas
simples ao longo do marcador: cada um passa a ter uma fita de cada
parental ao longo do trecho --- um heterodúplex com um mau pareamento
$A/a$. Se o \omterm{prop:b3:dna-repair:fidelity}{reparo de mau pareamento} converte o mau pareamento no
dúplex invadido para $A$, três cromátides carregam $A$ e uma carrega
$a$: $3{:}1$; convertido para $a$: $1{:}3$; deixado sem reparo, a
cromátide segrega $A$ e $a$ na primeira mitose depois da meiose, dando
dois esporos-filhos diferentes (segregação pós-meiótica, $5{:}3$ num
asco de oito esporos). O trecho de heterodúplex tem apenas de algumas
centenas a alguns milhares de pares de bases em torno da quebra
iniciadora, de modo que só marcadores dentro dele podem ser
convertidos, e eles estão, por construção, ao lado da permuta.
\end{solution}

\begin{solution}{pb:b3:dna-repair:1}
\textbf{1.} $10^{4} + 400 + 10^{4} \approx 2\times 10^{4}$ lesões por
dia, $7\times 10^{6}$ por ano --- grosso modo uma por quilobase do
genoma por ano.
\textbf{2.} $10^{-5}\times 10^{-2}\times 10^{-3} = 10^{-10}$ por bp;
$6.4\times 10^{9}\times 10^{-10} = 0.64$ mutações por divisão.
\textbf{3.} $0.64\times 0.015 \approx 0.01$ mutações codificadoras por
divisão; ao longo de $50$ divisões, cerca de $0.5$.
\textbf{4.} $10^{-7}$ por bp, $640$ mutações por divisão. Células com
deslizamento: sem reparo, $10^{-4}\times 40\times 10^{9} = 4\times 10^{6}$; com ele, $10^{-7}\times 40\times 10^{9} = 4000$ --- uma
diferença de mil vezes, visível como \omterm{def:b3:dna-repair:mmr}{instabilidade de microssatélites}.
\textbf{5.} A uracila é estranha ao DNA, é reconhecida pela
uracila-DNA-glicosilase e excisada: reparada. A timina vinda da
5-metilcitosina é uma base normal em frente a um G; nada a assinala
como errada fora da replicação, e ela se torna uma mutação C$\to$T ---
o mecanismo que esvaziou o genoma de \omterm{def:b3:chromatin-epigenetics:methylation}{CpG}.
\textbf{6.} A \omterm{def:b3:dna-repair:excision}{glicosilase} da 8-oxoG remove a base oxidada enquanto ela
ainda está em frente a C (antes da replicação); a segunda \omterm{def:b3:dna-repair:excision}{glicosilase}
remove um A mal incorporado em frente à 8-oxoG (depois da replicação,
restaurando uma chance de reparar o G); a hidrolase destrói o dGTP
oxidado para que ele não seja incorporado em frente a A. Quando as três
falham, a 8-oxoG pareia com A e a rodada seguinte fixa uma transversão
G$\cdot$C$\to$T$\cdot$A.
\textbf{7.} Normal $k = \ln 2/2 = \qty{0.35}{h^{-1}}$; XP $k = \ln 2/70
= \qty{0.0099}{h^{-1}}$.
\textbf{8.} Normal $10^{5} e^{-2.77} \approx 6200$; XP $10^{5}
e^{-0.079} \approx \num{92000}$.
\textbf{9.} Mutações $6.2$ e $92$ por célula; codificadoras $0.09$ e
$1.4$.
\textbf{10.} $500$ exposições: $47$ mutações codificadoras (normal) e
$690$ (XP), contra $0.5$ vindas da replicação --- a luz do sol domina a
carga de mutação da pele mesmo numa pessoa normal.
\textbf{11.} Sequência codificadora $9.6\times 10^{7}$ bp; os genes
condutores são $4\times 10^{4}/9.6\times 10^{7} = 4.2\times 10^{-4}$
dela. Média de acertos $690\times 4.2\times 10^{-4} = 0.29$;
$P(\ge 1) = 1 - e^{-0.29} =
0.25$. De $10^{9}$ células, $2.5\times 10^{8}$ carregam uma mutação condutora (criança normal: média $0.02$,
\qty{2}{\%}).
\textbf{12.} Um dímero não é uma mutação; ele só se torna uma quando
uma forquilha o copia por contorno propenso a erro, e por isso são as
células em divisão com muitos dímeros não reparados que mutam. A luz
ultravioleta alcança apenas a pele e o olho; os tecidos internos não
recebem dímeros e, no XP, estão quase normais.
\textbf{13.} $70$ quebras; após \qty{1}{h} só com NHEJ, $70\times
2^{-2} \approx 18$ focos.
\textbf{14.} $k_{\text{NHEJ}} = \qty{1.39}{h^{-1}}$, $k_{\text{HR}} =
\qty{0.17}{h^{-1}}$; fração da HR $0.17/1.56 = 0.11$.
\textbf{15.} Meia-vida $\ln 2/1.56 = \qty{27}{min}$; após \qty{1}{h},
$70\,e^{-1.56} \approx 15$ quebras.
\textbf{16.} $70\times 0.015\times 0.7 \approx 0.7$ genes rompidos por
célula.
\textbf{17.} $n = 70$: $2415$ pares $\times 3\times 10^{-4} = 0.72$
translocações. $n = 7$: $21\times 3\times 10^{-4} = 0.006$. O número de
quebras é $\propto D$, mas o número de \emph{pares} de quebras
simultâneas é $\propto n^{2} \propto D^{2}$.
\textbf{18.} $D = \qty{7}{h^{-1}}$, $L^{*} = 7/1.39 \approx 5$ quebras
presentes por vez: $10$ pares, $0.003$ translocações --- duzentas vezes
menos que a mesma dose num instante. O fracionamento deixa o tecido
normal reparar entre as frações e mantém poucas quebras simultâneas.
\textbf{19.} Reparo intacto: $e^{-2/1.5} = 0.26$ a \qty{2}{Gy},
$e^{-4} = 0.018$ a \qty{6}{Gy}. Sem NHEJ: $e^{-4} = 0.018$ e
$e^{-12} = 6\times 10^{-6}$. $D_{0}$ é a dose que deixa, em média, uma
\omterm{def:b3:dna-repair:lesions}{lesão} letal não reparada por célula (sobrevivência $1/e$).
\textbf{20.} Os \qty{11}{\%} de quebras de duas extremidades em G2 que
a HR teria reparado com acurácia passam agora pela NHEJ, e cada quebra
de forquilha de uma só extremidade --- que a NHEJ não consegue reparar
corretamente de modo algum --- é unida errada: ao longo dos anos o
genoma acumula deleções, translocações e mudanças de número de cópias,
as cicatrizes características dos tumores \emph{\omterm{def:b3:dna-repair:dsb}{BRCA}}.
\textbf{21.} $10^{4}\times 0.01 = 100$ quebras de uma só extremidade
por dia. Uma quebra de uma só extremidade não tem extremidade parceira;
a NHEJ só pode deixá-la como está ou uni-la a alguma extremidade sem
relação, produzindo uma translocação ou uma perda.
\textbf{22.} Um alelo \emph{BRCA2} funcional produz proteína suficiente
para a recombinação: as células normais reparam suas quebras de
forquilha e sobrevivem ao inibidor.
\textbf{23.} Uma segunda mudança de fase a jusante da primeira restaura
a fase de leitura e dá uma BRCA2 encurtada, mas em parte funcional, de
modo que a recombinação volta e o fármaco deixa de matar. Detecte
sequenciando o \emph{BRCA2} no tumor resistente ou no DNA tumoral
circulante, e funcionalmente pelo reaparecimento dos focos de \omterm{def:b3:dna-repair:dsb}{Rad51}
após a irradiação.
\textbf{24.} A morte celular precisa que o \omterm{prop:b3:dna-repair:fidelity}{reparo de mau pareamento}
engaje as bases alquiladas; um tumor de Lynch, que não o tem, é
tolerante --- o fármaco não consegue matá-lo (tais tumores são tratados
por outros meios, entre eles a imunoterapia, à qual sua pesada carga de
mutações os torna sensíveis).
\textbf{25.} Cerca de \num{92000} dímeros na forquilha da célula XP;
$10^{-10}$ mutações por par de bases por divisão; a HR repara
\qty{11}{\%} das quebras em G2.
\end{solution}
